\subsection{Painted errors and long paths}\label{sec:joint}

In one lane put into \(E\) the data replacement flags, of density \(2p\), and every site within a fixed multiple of \(H_n\) of any short or long bad center. Affected sites added by the latter operation will be called painted sites. Enlarge the paint multiplier once to accommodate the clipped comparisons in Section~\ref{sec:environment}. When helpful one may extend the dominating raw flag field to negative times or take a subset of the declared affected sites. In the probability arguments within a lane, products concern disjoint innovation supports. Counts across different lanes use unions.

\begin{lemma}[Joint sparse tests and finite-rank clusters]\label{lem:joint}
Fix the test multipliers used in the organization and cap arguments. The helper constants can be chosen so that for each rate in every \(I_n\) the sum of the failure probabilities for one prescribed \((j-1)\)-sparsity test in a fixed multiple of \(V^j\) at every rank \(j\ge1\) is less than \(1/4\), uniformly in their translations. For a fixed cap \(K\), mark each \(z\) for which some centered test, a fixed multiple of \(V^{i+1}\), fails \(i\)-sparsity for \(0\le i\le K\). The nearest- or diagonal-neighbor cluster of marked sites containing any specified site has an exponentially decaying max-norm radius tail. The exponent can depend on the rate and on \(K\). The same two conclusions hold for raw flags alone at all sufficiently small fixed rates, with a common choice of threshold for a given geometry.
\end{lemma}

\begin{proof}
Write \(b=\min(P,T)\), \(B_*=\max(P,T)\). For small \(p\), let \(j_s,j_l\) be the first ranks with \(b^{j_s}\ge p^{-17}\) and \(b^{j_l}\ge p^{-(b_h+6)}\). Fixed scale factors can be included. There are finite exponents \(d_s,d_l\) with \(B_*^{j_s}\le p^{-d_s}\), \(B_*^{j_l}\le p^{-d_l}\) after allowing the fixed scale factors. The first exponent depends only on the geometry, the second on the geometry and \(b_h\). At the respective scales the entire primitive support of a center event whose paint intersects a leaf fits in a bounded enlargement of the leaf. The leaf itself has side vector \(2wV^j\). Lemma~\ref{lem:sparse-witness} with \(u=2p\), after absorbing constants into a small exponent \(c>0\), bounds all raw nonsparsity probabilities by \(p^{c2^i}\), including bounded test enlargements.

Fix a nonsparse full box of rank \(i>j\). Iteration of Lemma~\ref{lem:sparse-witness} produces a binary tree with \(L=2^{i-j}\) disjoint leaf boxes of rank \(j\). There are at most \(C^{L-1}\) possible trees for fixed root. If each leaf event is measurable in a bounded enlargement of its leaf, the intersection graph of the enlarged leaves has degree at most a fixed \(\Delta_{\rm leaf}\): after scaling the two axes each neighboring original leaf is a disjoint unit rectangle lying within a bounded box. Select at least \(L/(\Delta_{\rm leaf}+1)\) disjoint supports in each deterministic tree. If the leaf probability is at most \(\alpha\), the product bound on the selected supports followed by a union over trees is at most
\begin{equation}\label{eq:tree}
(C\alpha^{1/(\Delta_{\rm leaf}+1)})^{2^{i-j}}.
\end{equation}

At \(j=j_s\) use as leaf event raw nonsparsity or a short bad center in the paint halo. The additional term is \(O((B_*^{j_s}+H_n)^2p^{M_s})\). Choose \(D_s>12d_s+10\) and \(M_s>\max(8d_s,(\Delta_{\rm leaf}+1)(D_s+1)+2d_s)\), then choose the short-event constants of Lemma~\ref{lem:stripes}. For small \(p\), \eqref{eq:tree} bounds raw-plus-short nonsparsity at \(i>j_s\) by \(p^{D_s 2^{i-j_s}}\). For \(i\le j_s\) we use raw nonsparsity plus the halo union bound directly. Now choose \(D_l>12d_l+10\) and \(M_l>\max(8d_l,(\Delta_{\rm leaf}+1)(D_l+1)+2d_l)\), then choose the channel count of Lemma~\ref{lem:stripes}. At \(j_l\) the raw-plus-short probability is smaller than any prescribed power of \(p\) as \(j_l-j_s\to\infty\). The long-center term is \(O((B_*^{j_l}+H_n)^2p^{M_l})\). Applying \eqref{eq:tree} again gives full nonsparsity at \(i>j_l\) bounded by \(p^{D_l2^{i-j_l}}\). A bounded cover transfers each of the estimates to each of the specified test shapes; strict slack in choosing the exponents absorbs the fixed factor. The eventual \(n_0\) absorbs every fixed prefactor. Both admissions use the chosen data values of \(P,T\), including \(P^n\) in the stripe period.

For the nested-test sum, bound the raw ranks through \(j_s\) by the raw double-exponential sequence, and charge short centers there by the sum bounded by \(O((j_s+1)(B_*^{j_s}+H_n)^2p^{M_s})\). Between \(j_s\) and \(j_l\), sum the raw-plus-short sequence. Charge long centers through \(j_l\) by \(O((j_l+1)(B_*^{j_l}+H_n)^2p^{M_l})\), and sum the full high-rank bounds above \(j_l\). All sums tend to zero as \(p\to0\) and are uniform in the test translations. This gives the stated threshold for the total.

For the cluster bound use three kinds of centered test marks: raw for \(i\le j_s\), raw-plus-short for \(j_s<i\le j_l\), and full for \(i>j_l\). The support radius of each is at most \(CB_*^i\) with a fixed constant, which includes the multiplier from \(V^{i+1}\). Also mark the short-center balls of radius \(CB_*^{j_s}\) and the long-center balls of radius \(CB_*^{j_l}\), whose probabilities are \(p^{M_s}\) and \(p^{M_l}\). With large enough fixed \(C\), every failed centered test belongs to the union of the marked balls, and their supports are contained in them. Indeed, for low ranks, one nearby painted center suffices for its own marked ball to cover the test center. In the middle ranks only long centers are charged that way, and all short events are included in the middle test. At high ranks the full test fits with its supports.

For each type \(\lambda\) let \(r_\lambda\ge1\) be its radius and \(a_\lambda\) its probability bound at any center. Our choices give, for all sufficiently small \(p\),
\begin{equation}\label{eq:ballsum}
\sum_\lambda C_1r_\lambda^4 \sqrt{a_\lambda}<\tfrac12 .
\end{equation}
For early ranks, the growth \(B_*^{4i}\) is dominated by \(p^{c2^i/2}\). In the middle range \(B_*^{4i}=B_*^{4j_s}B_*^{4(i-j_s)}\), where the first factor is at most \(p^{-4d_s}\); the choice of \(D_s\) gives summable decay including this factor. The late range uses \(d_l,D_l\) in the same computation. The choices \(M_s>8d_s\) and \(M_l>8d_l\) give the short- and long-ball terms. For raw flags alone this argument uses only the first type at all ranks.

A finite path from the root reaching max-distance \(t\) has a finite covering list of marked balls. Retain disjoint balls greedily in decreasing radius. Every discarded ball lies within the triple expansion of one retained ball. A chain of distinct triple-expanded retained balls connects the endpoints, allowing unit adjacency, and the sum of their original radii is at least \(t/8\). The event supports of the retained balls are disjoint. For an ordered list of types, the first center is at distance at most \(3r_1\) from the root, and successive centers can be at distance at most \(3r_i+3r_{i+1}+1\). The product of the corresponding two-dimensional lattice counts is bounded by \(\prod_i C_1r_i^4\), since each radius contributes twice to a quadratic count. This holds with an absolute geometric \(C_1\), which can be used in \eqref{eq:ballsum}. Through fixed \(K\) only finitely many types are needed, so \(\kappa=\min_\lambda[-\log a_\lambda/(2r_\lambda)]>0\). Reserve half the product probability for the factor \(e^{-\kappa t/8}\). Summation of the other half, including center counts and chain lengths, is bounded by the geometric series from \eqref{eq:ballsum}. This proves the tail bound and finishes the proof.
\end{proof}
